\section{Coomaraswamy on Evola’s Revolt}

The first translation of any of \textbf{Julius Evola}'s works into English was published in the prestigious Indian journal The Visva-Bharati Quarterly (Vol. V, Part IV, New Series, 1940), founded by the Nobel prize winning writer and poet, \textbf{Rabindranath Tagore}. \textbf{Ananda K Coomaraswamy} AKC wrote a brief introduction, which follows below, and the translation of the chapter ``Man and Woman'' from Revolt against the Modern World was made by his wife, Luisa Runstein. A few of Evola's essays appeared in the British journal East and West in the 1950s, but the full translation of Revolt did not appear in English until 1995.

For obvious reasons, Evola's book could receive more understanding in India than in the West. The modern deracinated Westerner is ignorant of his own past. Hence he can discern nothing recognizable in Evola's ``remarkable presentation and exposition of Traditional doctrine'' (AKC), which was characteristic not only of Indian civilization, but also of ancient Rome and the Middle Ages. Instead, he remains smug in the ``pretenses'' of the modern world.

AKC did point out a serious error due to Evola's misunderstanding of a text from the Aitareya Brahmana (VIII, 27, Page 341 of Arthur Keith's translation). In a Hindu wedding ceremony, the groom says to his bride:

\begin{quotex}
I am that, you are this, this you are, that I am; I am sky, you are earth.
\end{quotex}


When the King chooses his Purohita, i.e., the priest who offered sacrifices for the king since the gods would not accept those of the king directly, the purohita recites the groom's part to the king. Hence, in this regard, the priest is masculine in respect to the king. Evola assumed just the opposite, which led to the unfortunate characterization of the priesthood as lunar and regality as solar. Although this quotation was removed from subsequent editions of Revolt, this error persisted in Evola, even if inconsistently and ambiguously. This is AKC's note on that topic:

\begin{quotex}
We refer in particular to the citation of the conjugal \emph{rajnah purohita-varana-mantram} with which the priest takes the king as consort, saying, ``I am sky, you are earth''. Evola, inverting the roles of the king and priest, does so in a way that it was with these words that the king directed to the selected priest and consequently that the priest was to obey the king! The true relations were recognized by the first Sanskrit scholars (Oldenberg, for example), by Guenon in Spiritual Authority and Temporal Power, by Hocart in The Castes and were recently discussed in a very recent article on the various aspects of Hindu regality (Spiritual Authority and Temporal Power in the Indian Theory of Government, 1942).
\end{quotex}


AKC's introduction to the translation follows.

\begin{quotex}
Revolt against the Modern World, which contains the chapter titled ``Man and Woman'', whose translation follows, constitutes perhaps the most significant among the numerous works of the Italian philosopher Julius Evola. Evola's ``Revolt'' is, fundamentally, the affirmation of the universal metaphysical doctrine that has given form to traditional social institutions: thus this study treats social hierarchy (``caste''), the meaning of regality and empire, the relationship between the sexes, the mystery of rites, and from the same point of view interprets history and makes a destructive criticism of the pretenses of modern ``civilization''.

The book is not without defects: in particular he exalts beyond measure the importance of regality in relation to the priesthood (Kshatriya in relation to Brahma, raja in relation to Brahman), even erroneously interpreting Hindu texts to support his theses on this point. In addition, of everything amiss in a work of this type, in which the ``serenity'' of the heroic element results so rightly emphasized, he clearly shows an anti-Semitic prejudice, as can be seen in the superfluous characterization of Freud as a Jew. Nevertheless, this book constitutes a remarkable presentation and exposition of Traditional doctrine and could well serve as an introductory text for the student of anthropology and as a guide for Indology, especially for anyone who is interested in Hindu mythology and does not understand that, in the words of Evola, ``the passage from mythology to religion constitutes a humanist decadence.''

The chapter ``Man and Woman'' was chosen for translation because of its clear, uncompromising, and, we can add, dense peroration of the principles, that are reflected in the institutions and ideals, such as that of \emph{sati}, that are no longer comprehensible and that certainly are no longer held dear, not even as memories by our politicians and reformers who, ``whether by force or consent, were led to the accept the models of the West''.
\end{quotex}


That last quotation is from Peaks and Lamas, by Marco Pallis. AKC also provided ten footnotes to the translation, all but one of which supported the text with references to relevant passages in Hindu texts. The exception addressed Evola's claim that absolute ``dedication'' in love is typical for the woman in relation to the man. Giovanni Monastra describes AKC's objection:

\begin{quotex}
That, according to Coomaraswamy, is true in the limits in which the husband constitutes for the wife the ``symbol of God''. But, under the traditional aspect, in his turn the man also in relation to Divinity assumes a ``feminine'' role, rendering love with total dedication, without entailing any degradation, something different from what Evola means. Not even man can, therefore, be considered self-centered: the ``internal sufficiency'' which is mentioned in Revolt is a relative fact, not to be absolutized.
\end{quotex}

\paragraph{Traditional Mentality}

\begin{quotex}
We have touched upon only a very few of the ``motifs'' of folklore. The main point that we have wished to bring out is that the whole body of the motifs represent a consistent tissue of interrelated intellectual doctrines belonging to a \textbf{primordial wisdom} rather than to a primitive science; and that for this wisdom \emph{it would be almost impossible to conceive a popular, or even in any common sense of the term, a human origin}.
\flright{\textsc{Ananda Coomaraswamy}}
\end{quotex}


On a related note, I call attention to Coomaraswamy's perhaps misleadingly titled essay Primitive Mentality\footnote{\url{https://gornahoor.net/?page_id=4279}}. This version is missing the images from the original and most of the footnotes. I did include the most interesting footnotes, marked off in boxes.

Note that AKC refers to Evola twice in support of traditional society; this is certainly not a description of some fascist totalitarian dictatorship. Nevertheless, the traditional society is held together with a common mythology that orders their entire lives. Yet, within this structure, the individuals are free. When myth becomes religion, that is, it is compartmentalized into one aspect of life among others, then the culture begins to die.

AKC brings up the idea of a ``common mind'', without which an organic society is impossible. For the Traditional mind, based on transcendental principles, the idea of individual opinions is preposterous. Only in modern liberal societies is that conceivable. It does not make men free, but creates a atomized group of individuals who consider relationships as ``contracts'', not as living organic things. This is how the ``folk'' become the ``mass''.


\flrightit{Posted on 2012-05-28 by Cologero}

\begin{center}* * *\end{center}

\begin{footnotesize}
\begin{sffamily}
\texttt{Avery on 2012-05-28 at 19:44 said:}

Excellent sleuthing here. Evola's writing on Japan is similar; he intuits the entire metaphysical structure of the traditional Japanese nation, and indeed I am laboring to get his work translated over here, but sometimes his readings of specific elements are mistaken for lack of cultural context. I at first thought of Evola as a writer who is ``right even when he is wrong'', but I now realize that he can be wrong sometimes and still successfully communicate a powerful idea.

\hfill

\texttt{David on 2012-05-28 at 21:50 said:}

I read Evola's work extensivly, but never ''about'' him. This little text begs question then : where else did he misread or made errors ? Personnally, after reading him and myths (either Edda, Saint Graal, Gilgamesh, etc.) I seems to end to the same conclusion as he did : superiority of the Regal Divinity over both priest and warrior caste. Am I wronged ? If so, could someone tell me where to read something that would clarify all this ?

\hfill

\texttt{Sati Shankar on 2012-05-28 at 23:02 said:}

The male and female, king and priest,and the warrior, there is a mutual,(if this term can be used) overlap,one emerges from the other without any superiority of the either,... 'the Priest is the King himself, who appoints the Priest(Purohita)... .,without any numerical ordering... the opposing relation between as denoted is a conventional one,'not chronological or real order of coming into being', Fro example,the Son creates the Father as much as the Father the son... 'there can be no paternity without a filiation and vice versa,and that is what is meant by the `opposit relation'.The perspective of tradition has to emerge, as it does, from within which leads onto its whole.

\hfill

\texttt{Charlotte on 2012-05-30 at 13:23 said:}

Wow, mysterious kitchen talk about the planetary spheres, that's impressive... .and there was I thinking it mainly centred on bread and circuses :-)

\begin{verse}
`Now a voice so fair, ascending,\\
Fills the air with love unending,\\
Rises on the silver moonbeams\\
Woven from Apollo's sun streams.\\
`\,''Bold Orion, Starman leaping,\\
How my heart for you is beating.\\
I have set you there so thy fame\\
Lights the path of this, the sky-train!''\\
`Next she calls with gentle words\\
The creatures of her wooded world,\\
Speaks to them with tender charm\\
To keep the slightest safe from harm.\\
`\,``Sweet you are as honey, bee.\\
Bear and Stag, come follow me.\\
Jump with me across the river.''\\
Seeks she souls with bow and quiver.\\
Then the Goddess steps up on it --\\
Disc of night, the lamp of dreamers --\\
As the steeds with hooves of onyx\\
Take to flight with sweet Selene.
\end{verse}

\hfill

\texttt{FoaL on 2012-05-31 at 20:08 said:}
\begin{verse}
The Hero donned his solar face,\\
Departed from his holy place,\\
Armed and ready for a magic battle,\\
To defend and exalt his sacred cattle.\\
The phalanx of warriors will he lead,\\
The Gods to determine who shall bleed,\\
Ten thousand men all decked for war,\\
Rich in spirit, in money poor.\\
The women cry and halt their men,\\
Sayeth fair Charlotte ``We rather you remain in our love's den,''\\
``That leader of yours is too like Helios,''\\
``Aye'', affirmeth a wizened sage, ``I've said so many a time; too serious!''\\
But the fathers scoff at this hesitation,\\
``This world is for us procrastination!,''\\
``Your priority is procreation,''\\
``We are concerned with a higher calling,''\\
``Our enemies to see a-falling!''\\
In bursteth a champion, his name was HOO,\\
``Wenches, save your chatter for the kitchen or the loo.''\\
``On the morrow we commence our tour,''\\
``You are my lovers, but my destiny is Gornahoor.''\\
By night the chariots pass the wall,\\
After a day of feasting at the town hall.\\
The women gather to make solemn sacrifice.\\
Noble sun appeareth, come morning,\\
Cold as ice.
\end{verse}

\hfill

\texttt{Angolmois on 2012-06-01 at 04:50 said:}

``What is this supposed to mean?''

That he scorched himself in the solar plexus. The Sun burns and kills as much as it illuminates and gives light and life. Solar heroes have transcended all this to a supra-personal plane and that is why they shine without burning.

\hfill

\texttt{Sati Shankar on 2012-06-17 at 00:37 said:}

Friends, I hereby acknowledge all the participants of this thread that i will be referring to and using the information contained herein, in a Conference paper I am preparing. Due acknowledge will be cited in the paper also.

Regards,Sati Shankar

\hfill

\texttt{Michael M on 2021-05-29 at 10:00 said:}

``AKC brings up the idea of a ``common mind'', without which an organic society is impossible. For the Traditional mind, based on transcendental principles, the idea of individual opinions is preposterous.''

A common mind is indeed necessary for organic society. An interesting statement at first, I believe many would object to such a commonality between persons as somehow impinging on their individual freedoms. Again, this doesn't actually follow because the choice to be an idiot or evil while possible isn't precisely what someone wants for their lives as ``free'', especially when constructing or advocating for a society in which that ``freedom'' deters true freedom. Accepting wisdom isn't a punishment but a path to liberation and then all things flow from it, organically. Unfortunately, it does not appear we have hit ``rock bottom'' yet, accepting wisdom does indeed cause some pain especially in the beginning.

The idea of common mind has always struck a chord with the true essence of Catholic or what some thinkers were trying to put forward in Sobornost. Something common and universal to all that guides the people and gives boundaries to keep them together. The idea of such a common mind would be the representation of the sphere -- a vertical axis of principles that which all horizontal living revolves around with its own set of rules and regulations. Politics don't exist, because the principles cannot be debated, simply debated is the way to make decisions. But maybe I'm just romantically imagining life as a modern Beguine and Beghard in the countryside.

While societies cannot go back to such a mythology, the task at hands seems to incorporate the levels of mind common to all traditional societies. I don't mean synthesizing across societies, I mean appropriate to each society as the roots present are necessary for development. While the intellectual era of modernity has caused havoc, it is necessary to incorporate together. All functions of the human need to be acknowledged and then set to principles, I believe this is the task of the Christian, Christianizing all things with the Spirit. Maybe the virtue for the modern era is Chastity, which would be able to subdue the Lion (desires / consumption) and allow for Strength to come forth... Godspeed to everyone to convince others of such a virtue, let alone oneself.


\hfill

\end{sffamily}
\end{footnotesize}

